\documentclass[11pt]{article}

\usepackage[margin=1in]{geometry}
\usepackage[T1]{fontenc}
\usepackage{lmodern}
\usepackage{microtype}
\usepackage{amsmath,amssymb,amsthm,mathtools}
\usepackage[dvipsnames]{xcolor}
\usepackage[most]{tcolorbox}
\usepackage{tocloft}
\usepackage[hidelinks]{hyperref}

\definecolor{courseblue}{HTML}{1F4E79}
\definecolor{lightblue}{HTML}{EAF2F8}
\definecolor{lightgreen}{HTML}{EAF7EA}
\definecolor{lightpink}{HTML}{FCE8EF}
\definecolor{darkgray}{HTML}{333333}

\setlength{\parindent}{0pt}
\setlength{\parskip}{0.65em}
\renewcommand{\cftsecleader}{\cftdotfill{\cftdotsep}}
\renewcommand{\cftsubsecleader}{\cftdotfill{\cftdotsep}}

\theoremstyle{plain}
\newtheorem{theorem}{Theorem}[section]
\newtheorem{proposition}[theorem]{Proposition}
\newtheorem{lemma}[theorem]{Lemma}
\newtheorem{corollary}[theorem]{Corollary}
\theoremstyle{definition}
\newtheorem{definition}[theorem]{Definition}
\newtheorem{example}[theorem]{Example}
\newtheorem{exercise}[theorem]{Exercise}
\theoremstyle{remark}
\newtheorem*{remark}{Remark}

\tcolorboxenvironment{definition}{
  enhanced,
  colback=lightgreen,
  colframe=ForestGreen,
  boxrule=0.7pt,
  arc=2pt,
  left=6pt,
  right=6pt,
  top=5pt,
  bottom=5pt
}

\tcolorboxenvironment{exercise}{
  enhanced,
  colback=lightpink,
  colframe=WildStrawberry,
  boxrule=0.7pt,
  arc=2pt,
  left=6pt,
  right=6pt,
  top=5pt,
  bottom=5pt
}

\numberwithin{equation}{section}

% Numbered display.  (equation rather than the deprecated eqnarray;
% same output, and it respects amsmath spacing.)
\newcommand{\be}{\begin{equation}}
\newcommand{\ee}{\end{equation}}
\newcommand{\C}{\mathbb{C}}
\newcommand{\R}{\mathbb{R}}

\begin{document}

\begin{center}
  \colorbox{lightblue}{%
    \begin{minipage}{0.94\textwidth}
      \vspace{0.55em}
      \begin{tabular*}{\textwidth}{@{\extracolsep{\fill}}lr}
        {\color{courseblue}\bfseries\LARGE Lecture 3} &
        {\color{courseblue}\bfseries\LARGE MATH 411}\\[0.35em]
        \textbf{Fall 2026} & \textbf{Date: Sep 9, 2026}\\
        & \textbf{Instructor: Manish Patnaik}
      \end{tabular*}
      \vspace{0.45em}
    \end{minipage}}
\end{center}

\tableofcontents
\vspace{0.45em}

\section{Complex Differentiability}

Let $U\subseteq\C$ be open and let $f:U\to\C$.

\begin{definition}
We say that $f$ is \emph{complex differentiable} at $z_0\in U$ if the
limit
\be
  \lim_{\substack{h\to0\\ h\in\C\setminus\{0\}}}
  \frac{f(z_0+h)-f(z_0)}{h}
\ee
exists.  (Since $U$ is open, $z_0+h\in U$ for all sufficiently small
$h$, so the quotient makes sense.)  When the limit exists it is denoted
by $f'(z_0)$ and is called the \emph{derivative} of $f$ at $z_0$.

If $f$ is complex differentiable at every point of the open set $U$, we
say that $f$ is \emph{holomorphic} on $U$.
\end{definition}

The condition $h\to0$ is a limit in the complex plane: $h$ may approach
$0$ from any direction.  This makes complex differentiability much more
restrictive than ordinary one-variable differentiability, as the second
example below shows.

\begin{example}[The identity map]
For $f(z)=z$ and every $z_0\in\C$,
\[
  \frac{f(z_0+h)-f(z_0)}{h}
  =\frac{z_0+h-z_0}{h}=1.
\]
Thus $f$ is holomorphic on $\C$ and $f'(z)=1$.
\end{example}

\begin{example}[Complex conjugation]
Let $f(z)=\overline z$.  At any $z_0\in\C$,
\[
  \frac{f(z_0+h)-f(z_0)}{h}
  =\frac{\overline h}{h}.
\]
Writing $h=re^{i\theta}$ gives
\[
  \frac{\overline h}{h}=e^{-2i\theta}.
\]
The value depends only on the direction of approach, and it varies with
that direction: along the real axis ($\theta=0$) the quotient is
$1$, while along the imaginary axis ($\theta=\pi/2$) it is $-1$.  Hence
the limit does not exist, and complex conjugation is nowhere complex
differentiable.

Note that as a map of $\R^2$, $(x,y)\mapsto(x,-y)$ is linear, hence
as smooth as a function can be.  Complex differentiability is a
different, stronger, condition.  In Lecture~2 we computed for this
function that $u_x=1\ne-1=v_y$; we will see in Section~\ref{sec:CR} that
this is exactly the obstruction.
\end{example}

\begin{exercise}
Let $f(z)=|z|^2=z\overline z$.  Show that $f$ is complex differentiable
at $z_0=0$, with $f'(0)=0$, but at no other point of $\C$.  In particular
$f$ is not holomorphic on any open set, which is why the definition of
``holomorphic'' insists on differentiability throughout an open set.
\end{exercise}

\section{Basic Properties}

\begin{proposition}
If $f$ is complex differentiable at $z_0$, then $f$ is continuous at
$z_0$.
\end{proposition}

\begin{proof}
For $z\ne z_0$ we have
\be
  f(z)-f(z_0)
  =\frac{f(z)-f(z_0)}{z-z_0}\,(z-z_0).
\ee
As $z\to z_0$, the first factor tends to $f'(z_0)$ and the second tends
to $0$.  Therefore $f(z)\to f(z_0)$.
\end{proof}

\begin{proposition}[Differentiation rules]\label{prop:rules}
Suppose $f$ and $g$ are complex differentiable at $z_0$ and let
$c\in\C$.  Then $cf$, $f+g$ and $fg$ are complex differentiable at
$z_0$, and
\begin{align*}
  (cf)'(z_0)&=c\,f'(z_0),\\
  (f+g)'(z_0)&=f'(z_0)+g'(z_0),\\
  (fg)'(z_0)&=f'(z_0)g(z_0)+f(z_0)g'(z_0).
\end{align*}
If moreover $g(z_0)\ne0$, then $g\neq 0$ on a neighbourhood of $z_0$
(by continuity), the quotient $f/g$ is complex differentiable at $z_0$,
and
\be
  \left(\frac{f}{g}\right)'(z_0)
  =\frac{g(z_0)f'(z_0)-f(z_0)g'(z_0)}{g(z_0)^2}.
\ee
\end{proposition}

\begin{proof}
The proofs are word for word the same as for real functions of one
variable; they use only the algebra of limits and the continuity of $f$
and $g$ at $z_0$.
For example, try the proof of the product rule for yourself. (Hint: add and subtract
$f(z_0+h)g(z_0)$ in the numerator.)
\end{proof}

\begin{corollary}\label{cor:poly}
Every polynomial $p(z)=a_nz^n+\dots+a_1z+a_0$ is holomorphic on $\C$,
with
\be
  p'(z)=na_nz^{n-1}+\dots+2a_2z+a_1.
\ee
Every rational function $p/q$ is holomorphic on the open set
$\{z\in\C: q(z)\ne0\}$.
\end{corollary}

\begin{proof}
Constants and the identity map are holomorphic on $\C$ (with derivatives
$0$ and $1$), so by induction on $n$ and the product rule,
$(z^n)'=nz^{n-1}$.  The rest follows from
Proposition~\ref{prop:rules}.
\end{proof}

\begin{theorem}[Chain rule]\label{thm:chain}
Let $U,V\subseteq\C$ be open.  Suppose $f:U\to V$ is complex
differentiable at $z\in U$ and $g:V\to\C$ is complex differentiable at
$f(z)$.  Then $g\circ f$ is complex differentiable at $z$, and
\be
  (g\circ f)'(z)=g'(f(z))\,f'(z).
\ee
\end{theorem}

To prove the chain rule cleanly, it is useful to reformulate the
definition of differentiability.

\begin{lemma}[Linear approximation]\label{lem:linear}
A function $f$ is complex differentiable at $z$ if and only if there are
$\alpha\in\C$ and a function $\varphi$, defined on a neighbourhood of
$0$, such that
\be\label{eq:linear}
  f(z+h)-f(z)=\alpha h+h\,\varphi(h)
  \qquad\text{for all small $h$, including $h=0$,}
\ee
and $\varphi$ is continuous at $0$ with $\varphi(0)=0$.  In this case
$\alpha=f'(z)$.
\end{lemma}

\begin{proof}
If $f$ is differentiable at $z$, put $\alpha=f'(z)$, and for $h\ne0$
define
\be
  \varphi(h)=\frac{f(z+h)-f(z)}{h}-f'(z),
\ee
and set $\varphi(0)=0$.  Then~\eqref{eq:linear} holds for all small $h$
(trivially at $h=0$), and $\varphi(h)\to0=\varphi(0)$ by the definition
of $f'(z)$.  Conversely, dividing~\eqref{eq:linear} by $h\ne0$ and
letting $h\to0$ shows that the difference quotient tends to $\alpha$.
\end{proof}

\begin{proof}[Proof of the chain rule]
Put $w=f(z)$ and $k=f(z+h)-f(z)$, so that $f(z+h)=w+k$.  Applying the
lemma to $f$ at $z$ and to $g$ at $w$, we get functions $\varphi,\psi$,
each continuous at $0$ and vanishing there, with
\begin{align*}
  k&=f'(z)\,h+h\,\varphi(h),\\
  g(w+k)-g(w)&=g'(w)\,k+k\,\psi(k).
\end{align*}
The second identity holds for all small $k$, \emph{including $k=0$};
this matters, because $k$ can vanish for $h\ne0$ (for instance if $f$ is
constant).  Since $f$ is continuous at $z$, $k\to0$ as $h\to0$, and
hence $\psi(k)\to\psi(0)=0$.  Also $k/h\to f'(z)$.  Therefore, for
$h\ne0$,
\begin{align*}
 \frac{g(f(z+h))-g(f(z))}{h}
 &=g'(w)\,\frac{k}{h}+\frac{k}{h}\,\psi(k)\\
 &\longrightarrow g'(w)f'(z)+f'(z)\cdot0
 = g'(f(z))\,f'(z).
\end{align*}
This is the desired formula.
\end{proof}

\begin{remark}
The familiar ``proof'' from calculus, which writes
\[
  \frac{g(f(z+h))-g(f(z))}{h}
  =\frac{g(w+k)-g(w)}{k}\cdot\frac{k}{h},
\]
is not valid as written: the first factor is undefined whenever $k=0$.
The linear approximation lemma is precisely the device that avoids
dividing by $k$.
\end{remark}

\section{Real Total Differentiability}

Write $z=x+iy$ and
\be
  f(z)=u(x,y)+i\,v(x,y).
\ee
Identifying $\C$ with $\R^2$, the function $f$ corresponds to
\be
  F:U\subseteq\R^2\longrightarrow\R^2,
  \qquad
  F(x,y)=\begin{pmatrix}u(x,y)\\v(x,y)\end{pmatrix}.
\ee

\begin{definition}
The map $F$ is \emph{totally differentiable} at
$\mathbf a\in U$ if there is a real $2\times2$ matrix $D_{\mathbf a}$ such
that
\be
  \lim_{\mathbf h\to\mathbf0}
  \frac{\|F(\mathbf a+\mathbf h)-F(\mathbf a)-D_{\mathbf a}\mathbf h\|}
       {\|\mathbf h\|}=0.
\ee
The matrix $D_{\mathbf a}$ is unique and is called the derivative, or
Jacobian matrix, of $F$ at $\mathbf a$.
\end{definition}

If $F$ is totally differentiable at $\mathbf a$, then taking
$\mathbf h=t\mathbf e_1$ and $\mathbf h=t\mathbf e_2$ with $t\to0$ shows
that the first partial derivatives of $u$ and $v$ exist at $\mathbf a$
and that the derivative matrix is necessarily
\be
  D_{\mathbf a}=J_F(\mathbf a)
  =\begin{pmatrix}
      u_x&u_y\\
      v_x&v_y
    \end{pmatrix}_{\mathbf a}.
\ee
More generally, the directional derivative of $F$ at $\mathbf a$ in the
direction $\mathbf h$ is then $D_{\mathbf a}\mathbf h$.  The converse
fails: the existence of all directional derivatives at $\mathbf a$
does not by itself imply total differentiability at $\mathbf a$.

\begin{exercise}
Let $F:\R^2\to\R^2$ be given by $F(0,0)=(0,0)$ and
\[
  F(x,y)=\Bigl(\frac{xy^2}{x^2+y^4},\,0\Bigr)
  \qquad\text{for }(x,y)\ne(0,0).
\]
Show that every directional derivative of $F$ at the origin exists
(and that $J_F(0,0)$ is the zero matrix), but that $F$ is not even
continuous at the origin.  (Consider the curve $x=y^2$.)
\end{exercise}

\begin{remark}
A standard criterion from real analysis says that if the first partial
derivatives of $u$ and $v$ exist in a neighbourhood of $\mathbf a$ and
are continuous at $\mathbf a$, then $F$ is totally differentiable at
$\mathbf a$.
\end{remark}

\section{Relation with Real Analysis}\label{sec:CR}

Complex differentiability is real total differentiability with an extra
compatibility condition.  In the next lecture we will prove the following
precise statement.

\begin{theorem}\label{thm:CR}
Let $f=u+iv$ and let $F=(u,v)$.  Then $f$ is complex differentiable at
$z=x+iy$ if and only if $F$ is totally differentiable at $(x,y)$ and
\be
  \boxed{
    u_x=v_y,
    \qquad
    u_y=-v_x
  }
\ee
hold at $(x,y)$.  These are the \emph{Cauchy--Riemann equations}.  In
that case
\be
  f'(z)=u_x+i\,v_x=v_y-i\,u_y .
\ee
\end{theorem}

\begin{remark}[What the Cauchy--Riemann equations mean]
The Cauchy--Riemann equations say exactly that the Jacobian matrix has
the special form
\be
  J_F=\begin{pmatrix} a&-b\\ b&a\end{pmatrix},
  \qquad a=u_x,\quad b=v_x .
\ee
Under the identification $\C\cong\R^2$, multiplication by the complex
number $a+ib$ is the real-linear map with precisely this matrix.  So
Theorem~\ref{thm:CR} can be summarized as follows: \emph{$f$ is complex
differentiable at $z$ if and only if $F$ is totally differentiable at
$z$ and its derivative $D_z$ is not merely $\R$-linear but
$\C$-linear}, namely multiplication by $f'(z)$.  Geometrically, a
matrix of the form above with $(a,b)\ne(0,0)$ is a rotation composed
with a scaling by $\sqrt{a^2+b^2}=|f'(z)|$.  This is the source of the
angle-preserving (``conformal'') behaviour of holomorphic maps that we
will study later.
\end{remark}

\begin{exercise}
Using the partial derivatives computed in Lecture~2, check that
$f(z)=z^2$ satisfies the Cauchy--Riemann equations at every point and
that $f(z)=\overline z$ satisfies them nowhere.  Then compute $u,v$ for
$f(z)=|z|^2$ and show that the equations hold only at $z=0$.  Compare with the first section of
this lecture, and, for $z^2$, verify that
$u_x+i\,v_x$ agrees with the derivative $2z$ from
Corollary~\ref{cor:poly}.
\end{exercise}

\end{document}
